% CHU SO HUU: Ngọc Trâm | KHONG THEO PREAMBLE, KHONG BEGIN DOCUMENT TRONG THAN PHAN NAY.
\section{\texorpdfstring{\hspace{0.2em}}{}Phân tích phương sai một yếu tố và phương án Welch}
\label{sec:phuongphap-anova}

\subsection{Mục đích và lựa chọn biến}
ANOVA một yếu tố được dùng để đánh giá liệu trung bình của một đặc trưng định lượng có
giống nhau giữa Normal, DoS, Probe, R2L và U2R hay không. Bốn đặc trưng được chọn theo
ý nghĩa đo lường bổ sung nhau: \texttt{count} đại diện cường độ kết nối gần đây,
\texttt{serror\_rate} phản ánh lỗi SYN, \texttt{src\_bytes} đo lượng dữ liệu gửi từ
nguồn và \texttt{duration} đo thời gian duy trì kết nối. Chúng đồng thời đại diện cho
bốn dạng dữ liệu khác nhau (biến đếm, tỷ lệ bị chặn, dữ liệu đuôi phải và dữ liệu nhiều
số 0), đủ để trả lời các câu hỏi nghiên cứu có cơ sở mà không biến ANOVA thành bước sàng
lọc toàn bộ 41 đặc trưng. Các biến định tính như \texttt{service} hay \texttt{flag}
không phù hợp làm biến phụ thuộc của ANOVA.

\subsection{Giả thuyết}
Với mỗi đặc trưng, nhóm kiểm định
\[
H_0:\mu_{\mathrm{Normal}}=\mu_{\mathrm{DoS}}=\mu_{\mathrm{Probe}}
=\mu_{\mathrm{R2L}}=\mu_{\mathrm{U2R}},
\qquad
H_1:\text{có ít nhất một trung bình khác},
\]
ở mức ý nghĩa $\alpha=0{,}05$. Một kiểm định tổng thể có ý nghĩa chỉ cho biết ít nhất
một cặp khác nhau, không khẳng định mọi cặp đều khác.

\subsection{ANOVA cổ điển và kích thước ảnh hưởng}
Gọi $k$ là số nhóm, $N$ là tổng cỡ mẫu, $\bar x_i$ là trung bình nhóm và $\bar x$ là
trung bình chung. ANOVA cổ điển phân rã
\[
SS_B=\sum_{i=1}^{k}n_i(\bar x_i-\bar x)^2,\qquad
SS_W=\sum_{i=1}^{k}\sum_{j=1}^{n_i}(x_{ij}-\bar x_i)^2,
\]
và dùng
\[
F=\frac{\frac{SS_B}{k-1}}{\frac{SS_W}{N-k}}.
\]
Nhóm báo cáo thêm $\eta^2=\frac{SS_B}{SS_B+SS_W}$ và $\omega^2$ như các chỉ số mô
tả tỷ phần biến thiên gắn với nhóm. Khi kiểm định chính là Welch, các chỉ số dựa trên
tổng bình phương cổ điển được ghi rõ là mô tả, không suy ra trực tiếp từ Welch $F$.

\subsection{Phương án Welch}
Do cỡ mẫu rất mất cân bằng và phương sai khác nhau, kết quả chính dùng Welch ANOVA.
Đặt $w_i=\frac{n_i}{s_i^2}$, $W=\sum_iw_i$ và
$\bar x_w=\frac{\sum_iw_i\bar x_i}{W}$. Thống kê Welch được viết
\[
F_W=
\frac{
  \frac{\displaystyle\sum_iw_i(\bar x_i-\bar x_w)^2}{k-1}
}{
  \displaystyle 1+\frac{2(k-2)}{k^2-1}
  \sum_i\frac{\left(1-\frac{w_i}{W}\right)^2}{n_i-1}
},
\]
với $df_1=k-1$ và
\[
df_2=\frac{k^2-1}{
3\displaystyle\sum_i\frac{\left(1-\frac{w_i}{W}\right)^2}{n_i-1}
}.
\]
Nếu kiểm định tổng thể có ý nghĩa, Games--Howell được dùng cho 10 cặp vì không yêu cầu
phương sai hoặc cỡ mẫu bằng nhau. Hiệu chỉnh trong Games--Howell áp dụng cho các cặp
trong từng RQ; hiệu chỉnh Holm được áp dụng riêng cho bốn p-value tổng thể.

\subsection{Giả định, chẩn đoán và thang phân tích}
Các quan sát được giả định độc lập; phân bố trong nhóm không nên quá bất thường đối với
suy luận về trung bình; và ngoại lệ không được tự ý xóa để làm đẹp kết quả. Nhóm đánh
giá hình dạng bằng trung bình--trung vị, skewness, ECDF và Q--Q plot; đánh giá phương sai
bằng Brown--Forsythe; đồng thời xem xét chênh lệch cỡ mẫu và tỷ lệ bằng 0. Shapiro--Wilk
trên mẫu lớn chỉ là chẩn đoán bổ sung, không phải tiêu chí duy nhất để chọn kiểm định.
\texttt{count} và \texttt{serror\_rate} được kiểm định trên thang gốc để giữ câu hỏi
về trung bình gốc. Do lệch phải và cực trị mạnh, \texttt{src\_bytes} và
\texttt{duration} dùng $\ln(1+x)$; kết quả của chúng vì vậy nói về trung bình sau biến
đổi, không phải trực tiếp về số byte hay số giây trung bình trên thang gốc.

\subsection{Lấy mẫu phục vụ biểu đồ}
Mọi thống kê mô tả và kiểm định đều dùng toàn bộ 125.973 bản ghi. Chỉ bước trực quan hóa
dùng tối đa 5.000 quan sát/nhóm cho ECDF và 1.000 quan sát/nhóm cho Q--Q plot, với
\texttt{random\_state}=18015; nhóm nhỏ hơn ngưỡng được giữ toàn bộ. Các ngưỡng này đủ
lớn để thể hiện ổn định vị trí, đuôi và khối lượng tại 0, đồng thời tránh hàng chục nghìn
điểm chồng lấp, giảm kích thước hình và giữ nhãn có thể đọc khi in. Vì mẫu giảm chỉ dùng
để vẽ, nó không làm thay đổi $F$, p-value, khoảng tin cậy, effect size hay hậu kiểm.
Kết quả được trình bày tại Mục~\ref{sec:ketqua-anova}.

